% One Britannica scan as extracted by Nemotron Parse 3.5. Filled by src/nemotron_challenge/stages/latex.py.
% The TITLE, BODY and SOURCE markers (each wrapped in two percent signs) are replaced. Compile with XeTeX (tectonic).
\documentclass[10pt]{article}
\usepackage[a4paper,margin=18mm]{geometry}
\usepackage{fontspec}
\usepackage{multicol}
\usepackage{graphicx}
\usepackage[hidelinks]{hyperref}
\usepackage[normalem]{ulem}
\setlength{\parindent}{0pt}
\setlength{\parskip}{4pt}
\setlength{\columnsep}{8mm}
\begin{document}
\section*{Historical source: ENTOMOLOGY}
\begin{multicols}{2}
\subsection*{\textit{continued from the previous page}}
{\small\itshape [Margin: \emph{latro.} | \emph{argentata.} | fumigate. | \emph{saccata.} | \emph{extensa.} | \emph{ridulana.} | Aptera]}\par
a great height; to accomplish which, they climb up some eminence, and are walked about by the winds, filling the air with their threads. They are destroyed by the \sout{spikes} \textbf{sphinx} and ichneumon. Linnaeus has subdivided this genus according to the number and position of their \sout{eyes:} \textbf{eyes.} A. \sout{\_Eyes eight.\_ \_1.} \textbf{Eyes eight. 1.} Eyes placed \sout{thus\_,} \textbf{thus,} : \sout{fasnata.} \textbf{fasciata.} Of a silver \sout{color;} \textbf{colour;} the abdomen marked with yellowish \sout{bandes;} \textbf{bands;} the legs surrounded with brown rings. \sout{49.} \textbf{48.} A native of the island of \sout{Malachi} \textbf{Madeira.} * diadema Abdomen nearly \sout{globale} \textbf{globular} and reddish; brown, marked with a white cross composed of dots. 1. A native of Europe, on trees; it is a very beautiful insect. \_aquatica\_ \sout{Browny abbonen ovat} \textbf{Brown; abdomen oval} and of an ash colour; the back brown, marked with two dots. \sout{30.} \textbf{39.} A native of Europe, in stagnant waters, where it \sout{direct} \textbf{dives} to the bottom in search of \sout{the area.} \textbf{its food.} It takes up \sout{to} \textbf{its} winter \sout{quar-ters} \textbf{quarters} in an empty snail shell, the aperture of which it closes up with a web; jaws black, claws red. \sout{\#\#} 2. Eyes placed thus, : Thorax hairy, and of an ash colour; the abdomen oval, black, and spotted with red. 95. A native of America; large, the thorax oval, legs black, thighs pale. \_domestic.\_ \_obs.\_ \sout{\#\#} 3. Eyes placed thus, \_globosa.\_ Black; sides of the abdomen of a blood-red colour. \sout{99.} \textbf{69.} A native of Europe, in meadows. \sout{\#\#} 4. Eyes placed thus, \sout{Abdown} \textbf{Abdomen} white, and brown behind, and marked round the edge with six small projections. 70. A native of South America. \sout{\#\#} 5. Eyes placed thus, : \sout{Abdenom} \textbf{Abdomen} oval and brown, marked at the base with two white dots. 16. A native of Europe, in meadows; it \sout{washes} \textbf{watches} near the \sout{nest} \textbf{nests} of the \sout{larva} \textbf{larvee} of different insects, and \sout{make} \textbf{waits} their coming out, when it \sout{leaves} \textbf{seizes} on one after another, and sucks out their substance. \sout{Abdomoval,} \textbf{oval,} of a \sout{brownial, colour,} \textbf{brownish colour.} 40. A native of Europe, \sout{age} \textbf{in} gardens, and carries its eggs behind it enclosed in a bag; legs livid, marked with unequal brown rings set close together. \_Eyes placed thus, :\_ \sout{Abdemon} \textbf{Abdomen} long, of a silver colour tinged with green, the legs extended longitudinally. 22. A native of Europe, in woods; its legs are applied closely to the branches. \sout{\_Shining} \textbf{Shining} black; the abdomen hairy and black. \sout{73.\_} \textbf{73.} A native of America; it forms a nest under ground, and \sout{spin as} \textbf{spins a} large cylindrical web, which is covered with a \sout{id;} \textbf{lid;} its bite is very painful, and frequently occasions fever and delirium, which is soon removed by a

\subsection*{ENTOMOLOGY}
ENTOMOLOGY. gentle \sout{suborfics,} \textbf{sudorific ;} thorax marked with a large depressed circular \sout{spot; abdominal oval;} \textbf{spot ; abdomen oval ;} legs equal. \#\# 9. Eyes placed thus, : Black, thorax marked with a white line on the \sout{back\textbackslash{}*} \textbf{back.*} density 79. A native of Britain. Marked with triangular black spots on the back of \sout{torantsis,} \textbf{tarantula,} the abdomen; the legs spotted with black. 34. A native of the south of Europe, particularly Italy and Barbary. It is found in \sout{screens} \textbf{caverns} in \sout{agriculture} \textbf{argillaceous} soil; its bite was formerly supposed to be \sout{combble} \textbf{curable} by nothing but \sout{remains;} \textbf{music;} though it occasions a great deal of pain, it is almost never fatal. \sout{\#\#} 8. Eyes placed thus, Oval and oblong; thorax \sout{hair,} \textbf{hairy,} white; abdomen and \sout{\_pakkra\_.} \textbf{pulchra.} legs black, marked with yellow bands. 32. A native of Europe. \sout{\#\#} 9. Eyes placed thus, Black, \sout{abdominal} \textbf{abdomen} oval, forehead white, and leaps on gocuit, its prey. 84. It is to be found frequently in woods about Vienna; \sout{abdommed} \textbf{abdomen marked} with two compressed dots. \sout{\#\#} 10. Eyes placed thus, Thorax orbicular and convex, with a transverse \sout{cen-orientral tra} \textbf{central} excavation. 31. It is a native of South America, among \sout{others,} \textbf{trees,} where \sout{a} \textbf{it} preys upon the larger insects, and even small birds, dropping into their nests, and sucking their blood and \sout{legs;} \textbf{eggs;} it is of so enormous a size, that its \sout{flag} \textbf{fangs} may be compared to the \sout{latons} \textbf{talons} of a \sout{havk,} \textbf{hawk,} and its eyes are very large. \sout{\#\#} 11. Eyes placed thus, Of a grayish rusty colour, and slightly clouded; tho-truncets, rax globular, and slightly heart-shaped; abdomen some-what irregular, and marked with a faint longitudinal cross. 85. A native of Europe, on oak and other trees; abdommed marked with four impressed dots on the back; papilla three. \sout{\#\#} 12. Eyes placed thus, Abdomen oval, \sout{ostusely} \textbf{obtusely} conic behind, variegated \sout{comis.} \textbf{comas.} with brown and white; beneath black. 86. A native of Europe, among \sout{boshes;} \textbf{bushes;} it spins and spreads out a web, to the surface of which it attaches the insects it has \sout{emanated,} \textbf{ensnared,} after it has \sout{stacked} \textbf{sucked out} their \sout{joint. Tho-rax buck,} \textbf{juice. Thorax black,} slightly tinged with rusty colour, and nearly double behind; legs \sout{gravy,} \textbf{gray,} marked with brown rings; \sout{abdommed} \textbf{abdomen} marked \sout{as} \textbf{with} a \sout{whilish} \textbf{whitish} spot; beneath \sout{resembhing} \textbf{resembling} a horse shoe; \sout{papilla} \textbf{papillae} four. \sout{\#\#} 13. Eyes placed thus, Hairy; body oval, black, variegated with \sout{brown, albfiform.} \textbf{brown. albifrons.} 99. A native of Europe. B. Eyes six, placed thus, \sout{...} \textbf{.} Abdomen greenish, the sides yellow. 30. A \sout{naagents rive} \textbf{native} of Europe; in gardens.
\end{multicols}
\section*{Source record}
\begin{itemize}
\item Page ID: \texttt{britannica:309f765b1ecc83e0a7f7762e44b4a927a7555443:train:261}
\item Scan: \texttt{page-0261-881563669b04.jpg}, 600 x 776 px
\item Image SHA-256: \texttt{\small 881563669b04b487e86fd839d3c3ed37df89ebbf558afc819085fdba215bff61}
\item Extraction: \texttt{nvidia/nemotron-3.5-parse}, raw reply \texttt{artifacts/extract/raw/0261.json}
\item Text is Parse OCR output, normalized (long s as s, line-end hyphens joined). It is not corrected by hand.
\item \sout{Struck-through} words are Parse's reading; \textbf{bold} words are Super VL's correction, checked against the scan image. Both are model output.
\end{itemize}
\end{document}
